\documentclass[../relazione.tex]{subfiles}
\begin{document}

% ======================================================================
\section{Sistema di gestione (\protect\path{menu.py})}
% ======================================================================
Per semplificare l'esecuzione della pipeline, è stato implementato un sistema di gestione con due modalità operative:

\begin{enumerate}
    \item \textbf{Modalità automatica}: esecuzione sequenziale completa.
    \begin{lstlisting}
python menu.py [--update]
    \end{lstlisting}
    Esegue in sequenza tutti gli script (mapping, initialization, experiments, analyze, visualize) per tutti i modelli configurati. Il flag \texttt{--update} forza il ricalcolo di file esistenti;
    
    \item \textbf{Modalità interattiva}: menu testuale che permette l'interazione con l'utente.
    \begin{lstlisting}
python menu.py --menu [--update]
    \end{lstlisting}
    Presenta un menu che permette di:
    \begin{itemize}
        \item Eseguire singoli script della pipeline;
        \item Selezionare modelli specifici per l'analisi;
        \item Attivare la modalità update per singole operazioni.
    \end{itemize}
\end{enumerate}

\end{document}